\chapter{Observability and Incident Response}\label{ch:pl:observability-and-incident-response}

On 22 August 2013 the system that disseminates consolidated quotes for Nasdaq-listed stocks began to fail at 11:48 in the morning, with
intermittent outages on some of its channels. At 12:03 it stopped disseminating quotes altogether; a regulatory halt stopped trading in those
stocks on every venue; the technical fault was fixed by 12:24, and quoting and trading resumed at 15:25. A firm's monitoring that morning
could look perfect: every process up, every queue empty, every heartbeat on time --- because nothing was arriving. This chapter builds the
monitoring that would have noticed: signals about what users of the platform receive rather than about the processes that serve them,
objectives with \omterm{def:pl:observability-and-incident-response:budget}{error budgets}, alerts that page only when a person should act, and the incident practice that follows the page.

\section{Metrics, logs and traces}

\begin{definition}[Observability]\label{def:pl:observability-and-incident-response:observability}
\emph{Observability}\index{observability} is the degree to which the internal state of a running system can be inferred from what it
emits --- its \omterm{def:pl:observability-and-incident-response:metric}{metrics}, logs and traces --- without deploying new code to ask a new question.
\end{definition}

\begin{definition}[Metric, structured logging]\label{def:pl:observability-and-incident-response:metric}
A \emph{metric}\index{metric} is a named numeric measurement with labels, sampled or aggregated over time: a counter (messages received), a
gauge (queue depth), a histogram (latency, as in Book~13, chapter~5). \emph{Structured logging}\index{structured logging} writes each log event
as a record of named fields --- time, level, service, event, identifiers, values --- rather than a sentence, so that logs can be queried like
data.
\end{definition}

\begin{definition}[Distributed tracing, span]\label{def:pl:observability-and-incident-response:trace}
\emph{Distributed tracing}\index{distributed tracing} follows one unit of work --- a request, a market-data message, an order --- across the
services it passes through. Each step is a \emph{span}\index{span}, a unit of work with a start, an end, a trace identifier shared by all the
steps and the identifier of the step that caused it, so that the path can be rebuilt as a tree.
\end{definition}

The chapter's chain is the platform's own: capture (chapter~\ref{ch:pl:capturing-and-storing-tick-data}), the \omterm{def:pl:a-tick-store-on-open-formats:store}{tick store}, the \omterm{def:pl:position-and-pnl-services:service}{position
service} (chapter~\ref{ch:pl:position-and-pnl-services}) and real-time risk (chapter~\ref{ch:pl:real-time-risk}). \texttt{firm.observe} gives
each a \omterm{def:pl:observability-and-incident-response:metric}{metric} registry, structured log records and \omterm{def:pl:observability-and-incident-response:trace}{spans} in the OpenTelemetry sense, carried from one service to the next. A traced message
shows where its time goes (\cref{fig:pl:observability-and-incident-response:trace}): a fraction of a millisecond in capture, storage and
positions, then the wait in the risk service's queue --- 50~milliseconds on a normal day, a minute when the risk service falls behind.

\omcode{code/firm/observe/firm_observe.py}{103}{123}{Spans: a root span starts a trace, every child carries the trace's identifier and its
parent's, and a trace is rebuilt from them in time order.}\label{lst:pl:observability-and-incident-response:tracer}

\begin{omfigure}
\begin{tikzpicture}
\begin{semilogxaxis}[width=11cm,height=5.2cm,xbar,bar width=6pt,xmin=0.005,xmax=200000,y dir=reverse,ytick={0,1,2,3,4},
  yticklabels={capture,store,positions,risk queue,risk compute},xlabel={span duration (ms, log scale)},grid=major,grid style={gray!20},
  tick label style={font=\small},label style={font=\small},table/col sep=comma,log origin=infty,
  legend style={font=\small,at={(0.5,1.03)},anchor=south,legend columns=2,draw=none,
    /tikz/every even column/.append style={column sep=4mm}}]
\addplot[fill=omDef!40,draw=omDef] table[y=row,x=dur_ms]{figdata/platforms/28-observability-and-incident-response/trace_normal.csv};
\addplot[fill=omThm!50,draw=omThm] table[y=row,x=dur_ms]{figdata/platforms/28-observability-and-incident-response/trace_slow.csv};
\legend{normal day,slow risk consumer}
\end{semilogxaxis}
\end{tikzpicture}

\omcaption{The \omterm{def:pl:observability-and-incident-response:trace}{spans} of one traced market-data message through the chain (log scale). Capture, storage and positions take 0.02, 0.2 and
0.1~ms; the risk service computes in 1~ms. What changes when the risk service falls behind is its queue: 50~ms on a normal day, 60~seconds
at the height of the slow-consumer incident. Data: \texttt{fig\_observe.py}.}\label{fig:pl:observability-and-incident-response:trace}
\end{omfigure}

Each of the three answers a different question. \omterm{def:pl:observability-and-incident-response:metric}{Metrics} say that something is wrong, cheaply, for everything, all the time. Logs say what
happened in one service. Traces say where the time went across services, for the sampled units of work. Most trading firms have the first two
and not the third, and learn where a message spent its minute only by reading the logs of four services side by side, with four clocks.

\section{Service levels and error budgets}

\begin{definition}[Service-level indicator, service-level objective]\label{def:pl:observability-and-incident-response:sli}
A \emph{service-level indicator}\index{service-level indicator} (SLI) is a carefully defined measure of some aspect of the service a user
receives: here, the share of seconds in which the market data the risk service uses is at most five seconds old. A \emph{service-level
objective}\index{service-level objective} (SLO) is a target for an SLI over a period: here, 99.9\% of the trading seconds of a 30-day month.
\end{definition}

The indicator is chosen from the user's side. A heartbeat from the risk service says the process is alive; the age of its data says whether it
is doing its job. On the hook's morning every heartbeat was on time and every age grew by one second per second. An SLO differs from Book~14's
service-level agreement in whom it binds: an agreement is a promise to a customer with penalties, an objective is the firm's own target,
set tighter than any promise it makes.

\begin{definition}[Error budget]\label{def:pl:observability-and-incident-response:budget}
The \emph{error budget}\index{error budget} of an SLO is the amount of bad service it allows over its period --- one minus the target, times the
events --- which the service spends on incidents, maintenance and risky changes; when it is spent, reliability work takes priority over
features.
\end{definition}

A 99.9\% objective on 30 trading days of 23\,400 seconds allows 702 bad seconds a month. The chapter's quiet days --- the chain's normal
latency of tens of milliseconds, plus pauses (a burst, a collection, a slow disk) that arrive 0.8 times an hour and during which the age grows
one second a second --- spend 115 of them, 16\% of the budget, in 18 pauses of more than five seconds
(\cref{fig:pl:observability-and-incident-response:quiet}).

\begin{omfigure}
\begin{tikzpicture}
\begin{semilogyaxis}[width=11.5cm,height=5.2cm,xlabel={time of day},ylabel={oldest data age in the minute (s, log scale)},
  grid=major,grid style={gray!20},tick label style={font=\small},label style={font=\small},table/col sep=comma,
  xmin=0,xmax=390,ymin=0.03,ymax=200,xtick={0,90,180,270,360},xticklabels={09:30,11:00,12:30,14:00,15:30}]
\addplot[omDef,thin,const plot] table[x=minute,y=max_age]{figdata/platforms/28-observability-and-incident-response/quiet.csv};
\draw[omThm,dashed] (axis cs:0,5) -- (axis cs:390,5) node[pos=0.02,above right,font=\scriptsize]{5 s: the SLI's limit};
\draw[gray,dashed] (axis cs:0,30) -- (axis cs:390,30) node[pos=0.98,below left,font=\scriptsize]{30 s};
\draw[gray,dashed] (axis cs:0,60) -- (axis cs:390,60) node[pos=0.98,above left,font=\scriptsize]{60 s};
\end{semilogyaxis}
\end{tikzpicture}

\omcaption{A quiet day of the model: the oldest market-data age the risk service used in each minute. Outside the pauses the oldest age of a minute
stays about 0.1 to 0.3~seconds, a few times the chain's typical 50~ms; a few pauses of one or two seconds pass unnoticed; one pause at 14:35 lasts 56 seconds, crossing the SLI's five-second limit and every
static threshold below a minute. Data: \texttt{fig\_observe.py}.}\label{fig:pl:observability-and-incident-response:quiet}
\end{omfigure}

\section{Alerting that people act on}

A page wakes a person. It should mean that a person must act now, and it should come early enough for the action to matter. Static
thresholds on the data's age trade the two against each other; a rule on the \omterm{def:pl:observability-and-incident-response:budget}{error budget} does not have to.

\begin{definition}[Burn-rate alert]\label{def:pl:observability-and-incident-response:burn}
A \emph{burn-rate alert}\index{burn-rate alert} pages when the \omterm{def:pl:observability-and-incident-response:budget}{error budget} is being spent faster than a stated multiple of the rate that
would exactly exhaust it over the SLO's period, measured over a long window and confirmed over a short one: for a 99.9\% objective, Google's
SRE workbook recommends paging at 14.4 times over an hour (confirmed over five minutes) or 6 times over six hours (over thirty minutes).
\end{definition}

\omcode{code/firm/observe/firm_observe.py}{159}{176}{Multi-window burn-rate alerting: the share of bad seconds over each rule's long and
short windows, against the rate times the error budget's share.}\label{lst:pl:observability-and-incident-response:burn}

Three incidents are replayed on a quiet day at the hook's clock times: a feed stall from 12:03 lasting 21 minutes; intermittent stalls of 20
seconds every minute from 11:48 for a quarter of an hour; and a slow consumer from 14:00, a risk service whose lag grows by 0.05 seconds each
second. Five rules watch them (\cref{tab:pl:observability-and-incident-response:rules}).

\begin{table}[htbp]
\centering\small
\begin{tabular}{lrrrr}
\toprule
rule & feed stall & intermittent stalls & slow consumer & false pages a month\\
\midrule
process heartbeat & never & never & never & 0\\
data age above 5~s & 5~s & 5~s & 100~s & 18\\
data age above 30~s & 30~s & never & 600~s & 1\\
data age above 60~s & 60~s & never & 1\,200~s & 0\\
burn rate, 1~h and 6~h rules & 56~s & 191~s & 151~s & 1\\
\bottomrule
\end{tabular}
\caption{\omterm{def:pl:observability-and-incident-response:ttd}{Time to detect} each replayed incident, and pages over thirty quiet days, for five alerting rules. The heartbeat never fires; the
static thresholds buy speed with false pages or miss the intermittent stalls; the burn-rate rules detect all three within about three minutes
and page once a month on quiet days, for the day's 56-second pause.}
\label{tab:pl:observability-and-incident-response:rules}
\end{table}

\begin{definition}[Time to detect, time to restore]\label{def:pl:observability-and-incident-response:ttd}
The \emph{time to detect}\index{time to detect} of an incident is the time from its start to the first alert or report that a person acts on;
its \emph{time to restore}\index{time to restore} is the time from its start to the moment the service users receive is back within its
objective.
\end{definition}

\omterm{def:pl:observability-and-incident-response:ttd}{Time to restore} is measured on the service a user receives, from the incident's start; Book~14's mean time to repair is a property of a
component such as a circuit, averaged over its failures, from failure to repair.

The static thresholds show the trade-off (\cref{tab:pl:observability-and-incident-response:rules}). At five seconds they detect everything at
once and page 18 times a month on pauses that heal by themselves; within a few weeks nobody answers them. At thirty or sixty seconds the false
pages disappear, but the intermittent stalls --- the 11:48 phase of the hook --- never reach the threshold, and a slow consumer takes ten or
twenty minutes to be noticed. The burn-rate rules page 56 seconds into the stall, three minutes into the intermittent stalls and two and a half
into the slow consumer (\cref{fig:pl:observability-and-incident-response:slow}), with one page a month on quiet days: the 56-second pause, which
does spend a real share of the budget and deserves a look.

\begin{omfigure}
\begin{tikzpicture}
\begin{axis}[width=11.5cm,height=5.4cm,xlabel={minutes after the incident starts},ylabel={data age (s)},
  grid=major,grid style={gray!20},tick label style={font=\small},label style={font=\small},table/col sep=comma,
  xmin=-5,xmax=25,ymin=0,ymax=80,axis y line*=left]
\addplot[omDef,thick] table[x=t_min,y=age]{figdata/platforms/28-observability-and-incident-response/incident_slow.csv};
\draw[black,dashed] (axis cs:100/60,0) -- (axis cs:100/60,80) node[pos=0.97,left,font=\scriptsize]{5 s: 1.7 min};
\draw[omThm,thick] (axis cs:151/60,0) -- (axis cs:151/60,80) node[pos=0.85,right,font=\scriptsize]{burn rate: 2.5 min};
\draw[black,dashed] (axis cs:10,0) -- (axis cs:10,80) node[pos=0.97,left,font=\scriptsize]{30 s: 10 min};
\draw[black,dashed] (axis cs:20,0) -- (axis cs:20,80) node[pos=0.97,left,font=\scriptsize]{60 s: 20 min};
\node[omDef,font=\scriptsize,anchor=north west] at (axis cs:10.5,16) {age of the risk service's data};
\end{axis}
\end{tikzpicture}

\omcaption{The slow-consumer incident: from 14:00 the risk service processes messages slower than they arrive and the age of its data grows by
0.05 seconds each second. Its process stays alive throughout. The five-second threshold fires at 1.7 minutes, the burn-rate rule at 2.5, the
30- and 60-second thresholds after 10 and 20 minutes. Data: \texttt{fig\_observe.py}.}\label{fig:pl:observability-and-incident-response:slow}
\end{omfigure}

\begin{omfigure}
\begin{tikzpicture}
\begin{axis}[width=11.5cm,height=5.4cm,xlabel={minutes after 11:48},ylabel={burn rate (times the budget rate)},
  grid=major,grid style={gray!20},tick label style={font=\small},label style={font=\small},table/col sep=comma,
  xmin=-5,xmax=15,ymin=0,ymax=270]
\addplot[gray,thin] table[x=t_min,y=burn_5m]{figdata/platforms/28-observability-and-incident-response/incident_intermittent.csv};
\addplot[omThm,thick] table[x=t_min,y=burn_1h]{figdata/platforms/28-observability-and-incident-response/incident_intermittent.csv};
\draw[black,dashed] (axis cs:-5,14.4) -- (axis cs:15,14.4) node[pos=0.02,above right,font=\scriptsize]{14.4};
\draw[omThm,dotted,thick] (axis cs:191/60,0) -- (axis cs:191/60,270) node[pos=0.9,right,font=\scriptsize]{page: 3.2 min};
\node[gray,font=\scriptsize,anchor=north] at (axis cs:10,240) {5-minute window};
\node[omThm,font=\scriptsize,anchor=south] at (axis cs:11,66) {1-hour window};
\end{axis}
\end{tikzpicture}

\omcaption{The intermittent stalls: 20 seconds without data every minute from 11:48. The data age never exceeds 20 seconds, so thresholds at 30
and 60 seconds never fire; the burn rate over the last hour passes 14.4 --- one fiftieth of the month's budget in an hour --- after 3.2 minutes,
with the five-minute window confirming it. Data: \texttt{fig\_observe.py}.}\label{fig:pl:observability-and-incident-response:intermittent}
\end{omfigure}

\section{On call and incident command}

\begin{definition}[On-call rotation, runbook]\label{def:pl:observability-and-incident-response:oncall}
An \emph{on-call rotation}\index{on-call rotation} assigns, by a published schedule, the people who answer pages for a service at each hour,
with a secondary behind the primary. A \emph{runbook}\index{runbook} is the written procedure for a known alert or failure: what it means,
how to confirm it, what to do first, whom to call.
\end{definition}

\begin{definition}[Incident commander]\label{def:pl:observability-and-incident-response:ic}
The \emph{incident commander}\index{incident commander} holds the overall state of an incident and structures the response --- assigning
who changes the system, who communicates, who plans --- without doing the technical work themselves.
\end{definition}

In a trading firm the first minutes of an incident are about exposure, not diagnosis: is the firm trading on stale data, and should
strategies stop? The \omterm{def:pl:observability-and-incident-response:oncall}{runbook} for a data-age page therefore starts with the kill switches of Book~13 and the risk limits of Book~11, then the
diagnosis, and the \omterm{def:pl:observability-and-incident-response:ic}{incident commander}'s first question to the trading desk is what the firm holds. Only one group changes the system during an
incident, so that fixes do not collide; one person communicates, so that the desk, risk and compliance hear the same thing.

\begin{table}[htbp]
\centering\small
\begin{tabular}{lll}
\toprule
time & source & event\\
\midrule
12:03:00 & log & tickcap: last message on the quote channels\\
12:03:05 & alert & data age above 5~s: page\\
12:03:30 & alert & data age above 30~s: page\\
12:03:56 & alert & burn rate: page\\
12:04:00 & alert & data age above 60~s: page\\
12:24:00 & log & tickcap: quote channels resume\\
12:24:00 & log & rtrisk: data age back under 5~s\\
\bottomrule
\end{tabular}
\caption{The replayed feed stall's timeline, merged from alerts and structured logs by \texttt{firm.observe.timeline}. \omterm{def:pl:observability-and-incident-response:ttd}{Time to detect}:
56~seconds under the burn-rate rule; \omterm{def:pl:observability-and-incident-response:ttd}{time to restore}: 21 minutes.}\label{tab:pl:observability-and-incident-response:timeline}
\end{table}

\section{An incident replayed, and the review that follows}

\begin{definition}[Blameless review]\label{def:pl:observability-and-incident-response:blameless}
A \emph{blameless review}\index{blameless review} (post-incident review, postmortem) reconstructs an incident's timeline and contributing
causes without indicting any individual or team, on the assumption that everyone acted in good faith with the information they had, and
ends with actions that change the system rather than exhortations to be careful.
\end{definition}

The replayed stall's review writes itself from \cref{tab:pl:observability-and-incident-response:timeline}. The data stopped at 12:03; the
burn-rate page came at 12:03:56; the feed returned at 12:24 and the risk service's data was fresh again the same second. The questions it asks are
not who failed but what the system let happen: why the firm's dashboards showed green (they measured processes, not data); what the strategies
did with stale prices for the first minute (the kill switch's trigger was on process health, not data age); why the intermittent phase from
11:48 went unnoticed (thresholds above its twenty seconds). Its actions are changes: an SLI on data age for every consumer, burn-rate paging,
a kill switch on stale data, and the replay itself kept as a test (chapter~\ref{ch:pl:testing-and-continuous-delivery}) so that the next
change to the alerting is checked against it. An incident that becomes an operational loss event (Book~6, chapter~28) enters the firm's loss
data as well; the review is what makes the loss data useful.

\begin{dated}{2026-09}{The 2013 SIP outage and the SRE practices}\label{dat:pl:observability-and-incident-response:sip}
Nasdaq's vendor alert on the UTP SIP issue of 22 August 2013 gives the timeline: intermittent outages across multiple quote channels from
11:48, an outage on all quote channels at 12:03, the issue resolved and the environment stable at 12:24, connectivity re-established from
14:00, and all quoting and trading in Nasdaq-listed securities resumed at 15:25; the trade feeds were not affected, and a regulatory halt was
issued for all trading in those securities. Google's \emph{Site Reliability Engineering} defines SLIs, SLOs and the \omterm{def:pl:observability-and-incident-response:budget}{error budget} and describes
incident command and blameless postmortems; its workbook recommends, for a 99.9\% objective, paging on burn rates of 14.4 over one hour and 6
over six hours, each confirmed on a window one twelfth as long. OpenTelemetry defines a \omterm{def:pl:observability-and-incident-response:trace}{span} as a unit of work carrying trace, \omterm{def:pl:observability-and-incident-response:trace}{span} and parent
identifiers.
\end{dated}

\section{Tutorial: three hours of silence}\label{tut:pl:observability-and-incident-response:tut}

\begin{tutorial}
\textbf{Goal.} Instrument the chain, define its SLO, replay three incidents, compare alerting rules, and build the incident's timeline.
\textbf{End state:} \cref{tab:pl:observability-and-incident-response:rules,tab:pl:observability-and-incident-response:timeline}.
\begin{tutsteps}
\item \textbf{Traces}: \texttt{pl\_observe.traced(ts, lag)} on a normal day and during the slow consumer.
\item \textbf{Quiet days}: \texttt{quiet\_day(seed)} and \texttt{budget()}, the \omterm{def:pl:observability-and-incident-response:budget}{error budget} used in a month.
\item \textbf{Incidents}: \texttt{incident(kind)} for the three kinds.
\item \textbf{Rules}: \texttt{compare()}: \omterm{def:pl:observability-and-incident-response:ttd}{time to detect} and false pages for each rule.
\item \textbf{Timeline}: \texttt{stall\_timeline()} and the review.
\end{tutsteps}
\textbf{What to change next.} Add a ticket rule (burn rate 1 over three days) and count what it catches that the pages do not; make the
pauses twice as frequent and see which rules survive.
\end{tutorial}

\section{Build: observability}\label{bld:pl:observability-and-incident-response:observe}

\begin{build}
\bfield{Purpose} Know what the platform's users receive, page a person only when one must act, and reconstruct any incident from what the
platform emitted.
\bfield{Interface} \texttt{Registry} (\texttt{counter}, \texttt{gauge}, \texttt{histogram}), \texttt{log}, \texttt{Tracer} (\texttt{start},
\texttt{finish}, \texttt{trace}), \texttt{SLO}, \texttt{error\_budget\_used}, \texttt{threshold\_pages}, \texttt{burn\_rate\_pages},
\texttt{timeline}.
\bfield{Rules} SLIs measured from the user's side; every service emits \omterm{def:pl:observability-and-incident-response:metric}{metrics}, structured logs and \omterm{def:pl:observability-and-incident-response:trace}{spans} with propagated identifiers; pages
from burn rates on SLOs, tickets from slow burns; every page has a \omterm{def:pl:observability-and-incident-response:oncall}{runbook}; every incident has a commander and a \omterm{def:pl:observability-and-incident-response:blameless}{blameless review} with
changes.
\bfield{Acceptance tests} \texttt{code/firm/observe/tests/}: registry kinds and labels, a structured log line; a trace rebuilt from parent
links; budgets and burn rates; threshold pages with a minimum gap; a burn-rate page at the 52nd bad second of an outage and none before; a
merged timeline.
\bfield{Stretch} Exemplars linking a histogram bucket to a trace; tail-based trace sampling; alerts on the CUSUM of an SLI (Book~7,
chapter~13).
\end{build}

\begin{omsources}
\item Nasdaq, \emph{UTP Vendor Alert \#2013-9: UTP SIP Issue on Thursday, August 22, 2013}.
\item B.~Beyer, C.~Jones, J.~Petoff and N.~R.~Murphy (eds.), \emph{Site Reliability Engineering}, O'Reilly, 2016 (chapters 3, 4, 14, 15);
\emph{The Site Reliability Workbook}, 2018 (\emph{Alerting on SLOs}).
\item OpenTelemetry documentation, \emph{Traces}.
\item One Quant Book~13, chapters~5 and~24 (latency histograms, kill switches); Book~14, chapter~15 (service-level agreements, mean time to
repair).
\end{omsources}

\section{Exercises}

\begin{exercise}[$\star$]\label{exo:pl:observability-and-incident-response:1}
How many bad seconds does a 99.9\% objective allow over thirty days of 23\,400 seconds, and how many did the quiet month spend?
\end{exercise}

\begin{exercise}[$\star$]\label{exo:pl:observability-and-incident-response:2}
Why does the process heartbeat never fire in any of the three incidents?
\end{exercise}

\begin{exercise}[$\star$]\label{exo:pl:observability-and-incident-response:3}
How long does the burn-rate rule take to page on a complete stall, and where does the number come from?
\end{exercise}

\begin{exercise}[$\star\star$]\label{exo:pl:observability-and-incident-response:4}
Why do the thirty- and sixty-second thresholds never see the intermittent stalls?
\end{exercise}

\begin{exercise}[$\star\star$]\label{exo:pl:observability-and-incident-response:5}
What does the short window of a burn-rate rule add to the long one?
\end{exercise}

\begin{exercise}[$\star\star$]\label{exo:pl:observability-and-incident-response:6}
What is the replayed stall's \omterm{def:pl:observability-and-incident-response:ttd}{time to restore}, and why is it not a mean time to repair?
\end{exercise}

\begin{exercise}[$\star\star\star$]\label{exo:pl:observability-and-incident-response:7}
\emph{Coding.} Make the pauses twice as frequent (\texttt{pauses\_per\_hour=1.6}) and recount the false pages of the five-second threshold and
of the burn-rate rules over thirty days.
\end{exercise}

\begin{exercise}[$\star\star\star$]\label{exo:pl:observability-and-incident-response:8}
\emph{Find the flaw.} \textquotedbl{}Our monitoring was fine: every service reported healthy throughout the outage.\textquotedbl{}
\end{exercise}

\section{Problem: Three Hours of Silence}

\begin{problem}[{Weekend problem --- three hours of silence}]\label{pb:pl:observability-and-incident-response:1}
The chapter's chain, SLO, quiet days, incidents and rules.

\medskip\noindent\textbf{Part I --- Signals.}
\begin{enumerate}
\item What do \omterm{def:pl:observability-and-incident-response:metric}{metrics}, logs and traces each answer?
\item Where does a traced message spend its time on a normal day and during the slow consumer?
\item Why is the SLI the data's age and not the process's health?
\item What does the SLO allow, and what do quiet days spend?
\item How does an SLO differ from a service-level agreement?
\end{enumerate}

\medskip\noindent\textbf{Part II --- Rules.}
\begin{enumerate}[resume]
\item What are the three replayed incidents?
\item What does each static threshold detect, and how fast?
\item How many false pages a month does each rule send?
\item What do the burn-rate rules detect, and how fast?
\item Why is the burn-rate rule's one false page a month acceptable?
\end{enumerate}

\medskip\noindent\textbf{Part III --- Response.}
\begin{enumerate}[resume]
\item What does the \omterm{def:pl:observability-and-incident-response:ic}{incident commander} do and not do?
\item What comes first in the \omterm{def:pl:observability-and-incident-response:oncall}{runbook} for a data-age page in a trading firm?
\item What is the replayed stall's timeline?
\item What are its \omterm{def:pl:observability-and-incident-response:ttd}{time to detect} and \omterm{def:pl:observability-and-incident-response:ttd}{time to restore}?
\item What does a \omterm{def:pl:observability-and-incident-response:blameless}{blameless review} produce?
\end{enumerate}

\medskip\noindent\textbf{Part IV --- The verdict.}
\begin{enumerate}[resume]
\item State the \emph{named result}: the \omterm{def:pl:observability-and-incident-response:ttd}{time to detect} a feed stall and a slow consumer under a static threshold and under \omterm{def:pl:observability-and-incident-response:burn}{burn-rate
alerts}, and the false pages a month each costs, on the replayed incident and on thirty quiet days.
\item Which rule would you deploy, and with what ticket rule behind it?
\item What would have made the firm's dashboards turn red on the hook's morning?
\item What changes would the review of the replayed stall require?
\item In one sentence: what should a page mean?
\end{enumerate}
\end{problem}

\section{Interview questions}

\begin{interviewq}[$\star$ \iqroles{developer}]\label{iq:pl:observability-and-incident-response:1}
What is the difference between monitoring a process and monitoring a service?
\end{interviewq}

\begin{interviewq}[$\star\star$ \iqroles{developer}]\label{iq:pl:observability-and-incident-response:2}
What are an SLI, an SLO and an \omterm{def:pl:observability-and-incident-response:budget}{error budget}? Give one for a market-data platform.
\end{interviewq}

\begin{interviewq}[$\star\star$ \iqroles{developer}]\label{iq:pl:observability-and-incident-response:3}
Your team gets thirty pages a week and ignores most of them. What do you change?
\end{interviewq}

\begin{interviewq}[$\star\star$ \iqroles{developer}]\label{iq:pl:observability-and-incident-response:4}
A message takes a minute to go from the feed to the risk system. How do you find out where?
\end{interviewq}

\begin{interviewq}[$\star\star\star$ \iqroles{developer}]\label{iq:pl:observability-and-incident-response:5}
Design the \omterm{def:pl:observability-and-incident-response:observability}{observability} and incident process for a trading firm's data and risk platform.
\end{interviewq}
